\documentclass[10pt,a4paper]{amsart}
\usepackage[margin=19mm]{geometry}
\usepackage{amsmath,amssymb,mathtools}
\usepackage[scr=rsfs,cal=euler]{mathalfa}
\usepackage{xfrac}
\usepackage[unicode,hidelinks,bookmarks=false]{hyperref}
\newcommand{\ie}{i.e.\ }
\newcommand{\cf}{cf.\ }
\newcommand{\resp}{respectively\ }
\newcommand{\Subsection}{\subsection}
\providecommand{\nobreakdash}{\nobreak\hbox{-}}
\setlength{\emergencystretch}{2em}
\multlinegap0pt
\usepackage{amssymb}
\newtheorem{theorem}{Theorem}
\newtheorem{corollary}{Corollary}
\newtheorem{lemma}{Lemma}
\newtheorem{remark}{Remark}
\newtheorem{claim}{Claim}
\title{Spectra of generators of hyperbolic composition and weighted composition semigroups}
\author{Yong-Xin Gao, Ze-Hua Zhou}
\date{Source: arXiv:2605.22048v1 (CC BY 4.0)}
\begin{document}
\maketitle

\noindent\emph{Source and licence.} Spectra of generators of hyperbolic composition and weighted composition semigroups. Version arXiv:2605.22048v1 (CC BY 4.0), distributed by arXiv under the Creative Commons Attribution 4.0 International licence, http://creativecommons.org/licenses/by/4.0/, as recorded in the arXiv licence metadata for this exact version and shown on the abstract page https://arxiv.org/abs/2605.22048v1. Full licence notice: \texttt{licenses/arxiv-2605.22048-CC-BY-4.0.txt}.\par\smallskip

\noindent\textit{Editorial scope.} Counted unit: the article's theorem main (source0.tex lines 184--197). The article's complete original proof of it is delivered with it (source0.tex lines 889--957, one \texttt{proof} environment of 69 lines, which the article places away from the statement). No mathematics is written, repaired or re-derived: the statement and the proof are the article's own lines, in the article's own notation, and no retained line is substituted or re-flowed.  Retained uncounted as the source-local context the proof needs: lines 170--172; lines 433--436; lines 482--485; lines 514--516; lines 529--538; lines 606--608; lines 661--663; lines 665--667; lines 674--678; lines 702--706; lines 728--732; lines 794--802; lines 847--851.  Nothing was omitted from any retained span, so the package carries no omission record.  No retained line contains a citation, so the wrapper carries no bibliography and no bibliography entry.  No retained line contains a figure environment, an image-inclusion command or drawing code, so no image is delivered and the package declares no asset dependency.  Editorial formatting only, in addition: an AMS \texttt{amsart} preamble that restates the article's own declarations and macros used by the retained lines, namely the environments \texttt{theorem}, \texttt{corollary}, \texttt{lemma}, \texttt{remark}, \texttt{claim} on separate counters, together with the standard packages those lines need.  The retained environments print in this document as Theorem 1, Corollary 1, Lemma 1, Remark 1, Claim 1 in order of appearance, whereas the article prints the same items with the numbering of its own full document.  The complete native source of the article as distributed in the arXiv e-print for this exact version is bundled unchanged as \texttt{sources/201107/source0.tex}.
\begin{align}\label{gamma}
\gamma_j=2\alpha_j/p+ \operatorname{Re}\beta_j,
\end{align}

\begin{theorem}\label{rad}
Let $(u_tC_{\varphi_t})$ be a semigroup of bounded operators on $A^p$, where $(\varphi_t)$ is a hyperbolic semiflow. Suppose that $(\varphi_t)$ is analytic at its self-contact points, and $(u_t)$ is continuous at these points. Then for any $t\geq0$ the spectral radius of $u_tC_{\varphi_t}$ satisfies
$$r(u_tC_{\varphi_t})\leq\max\bigl\{\mathrm{e}^{\gamma_j t} : j\geq0\bigr\}.$$
\end{theorem}

\begin{corollary}\label{spectralbound}
Let $(u_tC_{\varphi_t})$ be a $C_0$-semigroup on $A^p$, where $(\varphi_t)$ is a hyperbolic semiflow. Suppose $(\varphi_t)$ is analytic at its self-contact points, and $(u_t)$ is continuous at these points. Then the spectrum of its generator $A$ satisfies
$$\sigma(A)\subset\left\lfloor-\infty, \max_{j\geq 0}\gamma_j\right\rceil.$$
\end{corollary}

\begin{align}\label{difsolu}
F(z)=\frac{ \mathrm{e}^{\lambda h(z)}}{v(z)}\left(K-\int_0^z \mathrm{e}^{-\lambda h(\xi)}h'(\xi)v(\xi)f(\xi) \mathrm{d}\xi\right),
\end{align}

\begin{theorem}\label{eigenvalue}
Let $(u_tC_{\varphi_t})$ be a $C_0$-semigroup on $A^p$, where $(\varphi_t)$ is a hyperbolic semiflow. Suppose $(\varphi_t)$ is analytic at its self-contact points, and $(u_t)$ is continuous at these points.
\begin{enumerate}
\item If $\gamma_1<\gamma_0$, then 
$$\overline{\sigma_p(A)}=\lfloor\gamma_1,\gamma_0\rceil,$$
and each interior point of this strip  (or half-plane) is an eigenvalue of $A$.

\item  If $\gamma_1>\gamma_0$, then $\sigma_p(A)=\emptyset$.
\end{enumerate}
\end{theorem}

\begin{lemma}\label{jifenshoulian}
If $ \operatorname{Re}\lambda>\gamma_0$, then for each $f\in A^p$, the integral $\int_z^{\zeta_0}\omega_{\lambda,f}$ gives an analytic function with respect to $z\in\mathbb{D}$.
\end{lemma}

\begin{lemma}\label{zaibianyige}
If $\operatorname{Re}\lambda<\gamma_j$ for some $j\geq 1$, then $\int^z_{\zeta_j}\omega_{\lambda,f}$ gives an analytic function with respect to $z\in\mathbb{D}$ for each $f\in A^p$.
\end{lemma}

\begin{remark}\label{zuihouremark}
The proofs of Lemma \ref{jifenshoulian} and \ref{zaibianyige} shows that when considering the integrals  $\int^z_{\zeta_j}\omega_{\lambda,f}$ for $j\geq0$, the integration path can be taken as any smooth curve from $z$ to $\zeta_j$ that approaches $\zeta_j$ nontangentially.
\end{remark}

\begin{claim}\label{>0}
If $ \operatorname{Re}\lambda>\gamma_0$, then for any $f\in A^p$,
$$\frac{ \mathrm{e}^{\lambda h(z)}}{v(z)}\left(K-\int_0^z\omega_{\lambda,f}\right)\in A^p_{\zeta_0}$$
if and only if $K=\int_0^{\zeta_0}\omega_{\lambda,f}$.
\end{claim}

\begin{claim}\label{others}
Let $\zeta$ be a point on $\partial \mathbb{D}$ that is not a fixed point of $(\varphi_t)$. Then
$$\frac{ \mathrm{e}^{\lambda h(z)}}{v(z)}\left(K-\int_0^{z}\omega_{\lambda,f}\right)\in A^p_{\zeta}$$
for every $f\in A^p$ and $\lambda, K\in\mathbb{C}$. 
\end{claim}

\begin{claim}\label{>j}
If $ \operatorname{Re}\lambda>\gamma_j$ for some $j\geq 1$, then
$$\frac{ \mathrm{e}^{\lambda h(z)}}{v(z)}\left(K-\int_0^{z}\omega_{\lambda,f}\right)\in A^p_{\zeta_j}$$
for any $f\in A^p$ and $K\in\mathbb{C}$. 
\end{claim}

\begin{claim}\label{<j}
If $ \operatorname{Re}\lambda<\gamma_j$ for some $j\geq 1$, then for any $f\in A^p$,
$$\frac{ \mathrm{e}^{\lambda h(z)}}{v(z)}\left(K-\int_0^z\omega_{\lambda,f}\right)\notin A^p_{\zeta_j}$$
whenever $K\ne\int_0^{\zeta_j}\omega_{\lambda,f}$.

Moreover, if we also have $\beta_0\ne-\infty$, then 
$$\frac{ \mathrm{e}^{\lambda h(z)}}{v(z)}\int^z_{\zeta_j}\omega_{\lambda,f}\in A^p_{\zeta_j}.$$

\end{claim}

\begin{claim}\label{<0}
Suppose $\gamma_1\ne-\infty$. If $ \operatorname{Re}\lambda<\gamma_0$, then 
$$\frac{ \mathrm{e}^{\lambda h(z)}}{v(z)}\left(K-\int_0^{z}\omega_{\lambda,f}\right)\in A^p_{\zeta_0}$$
for any $f\in A^p$ and $K\in\mathbb{C}$. 
\end{claim}

\begin{theorem}\label{main}
Let $(u_tC_{\varphi_t})$ be a $C_0$-semigroup on $A^p$, where $(\varphi_t)$ is a hyperbolic semiflow. Suppose $(\varphi_t)$ is analytic at its self-contact points, and $(u_t)$ is continuous at these points. Then the spectrum of the generator $A$ of $(u_tC_{\varphi_t})$ is given as follows:
\begin{enumerate}
\item If $\gamma_0 \geq \gamma_1$, then
$$\sigma(A) = \lfloor-\infty, \gamma_2\rceil\cup\lfloor\gamma_1,\gamma_0\rceil.$$

\item If $\gamma_2 < \gamma_0 < \gamma_1$, then
$$\sigma(A) = \lfloor-\infty,  \gamma_2\rceil\cup\lfloor\gamma_0 ,\gamma_1\rceil.$$

\item If $\gamma_0 \leq \gamma_2$, then
$$\sigma(A) = \lfloor-\infty, \gamma_1\rceil.$$
\end{enumerate}
Here the numbers $\gamma_j$ are defined in \eqref{gamma} and are decreasing for $j\geq 1$.
\end{theorem}

\begin{proof}[\textbf{Proof of Theorem \ref{main}}]
We rearrange (i)-(iii) of Theorem \ref{main} into the following five cases:

\noindent\emph{ Case \uppercase\expandafter{\romannumeral1}: $\gamma_1=-\infty$.} By Corollary \ref{rad} and Theorem \ref{eigenvalue}, we have
$$\sigma(A)=\lfloor-\infty, \gamma_0\rceil$$
in this case.



\noindent\emph{ Case \uppercase\expandafter{\romannumeral2}: $\gamma_0>\gamma_1>-\infty$.} Again by Theorem \ref{eigenvalue}, the strip $\lfloor\gamma_1, \gamma_0\rceil$ is contained in $\sigma(A).$

Suppose $\gamma_2<\operatorname{Re}\lambda<\gamma_1$. Claim \ref{>j}, \ref{<j}, and \ref{<0} show that for any $f\in A^p$,
$$\frac{ \mathrm{e}^{\lambda h(z)}}{v(z)}\int_z^{\zeta_1} \omega_{\lambda,f}\in A^p.$$
Hence, by taking $K=\int_0^{\zeta_1} \omega_{\lambda,f}$ in \eqref{difsolu}, we see that $\lambda-A$ is surjective, so $\lambda\notin\sigma(A)$.

Now suppose $\operatorname{Re}\lambda<\gamma_2$. Fix an arbitrary $f\in A^p$. If we assume that there existed a constant $K$ such that
$$\frac{ \mathrm{e}^{\lambda h(z)}}{v(z)}\left(K-\int_0^z \omega_{\lambda,f}\right)\in A^p,$$
then applying Claim \ref{<j} at $\zeta_1$, we see that the only possible choice of $K$ is $K=\int_0^{\zeta_1} \omega_{\lambda,f}$. But if we consider at  $\zeta_2$, then for the same reason we must have $K=\int_0^{\zeta_2} \omega_{\lambda,f}.$ Therefore, we have 
\begin{align}\label{linghua}
\int_{\zeta_1}^{\zeta_2} \omega_{\lambda,f}=0
\end{align}
for every $f$ in the range of $\lambda-A$. However, equation \eqref{linghua} cannot hold for all $f\in A^p$. In fact,  since $A^p$ is invariant under automorphisms, we may assume, after composing with a suitable automorphism, that $\zeta_1=-1$ and $\zeta_2=1$. Then by Remark \ref{zuihouremark}, the integral in \eqref{linghua} can be taken along the real segment $[-1,1]$. So equation \eqref{linghua} gives
$$\int_{-1}^{1} \mathrm{e}^{-\lambda h(x)}v(x)h'(x)f(x)\mathrm{d}x=0,$$
for all $f\in A^p$. By taking $f$ as polynomials, we have $\mathrm{e}^{-\lambda h(x)}v(x)h'(x)$ is identity zero on $(-1,1)$, which is  impossible. So $\lambda-A$ is not surjective, and therefore $\lfloor-\infty, \gamma_2\rceil\subset\sigma(A)$. To sum up, we have 
$$\sigma(A)=\lfloor-\infty, \gamma_2\rceil\cup\lfloor\gamma_1, \gamma_0\rceil$$
in this case.







\noindent\emph{ Case \uppercase\expandafter{\romannumeral3}: $\gamma_2<\gamma_0<\gamma_1$.} As shown in Case \uppercase\expandafter{\romannumeral2}, we have $\lambda\in\sigma(A)$ when $\operatorname{Re}\lambda\leq\gamma_2$, and $\lambda\notin\sigma(A)$ when $\gamma_2<\operatorname{Re}\lambda<\gamma_0$.

Now suppose $\gamma_0<\operatorname{Re}\lambda<\gamma_1$. Assume that $f$ belongs to the range of $\lambda-A$. Then according to Claim \ref{>0} and \ref{<j}, we must have
$$K=\int_0^{\zeta_0} \omega_{\lambda,f}=\int_0^{\zeta_1} \omega_{\lambda,f}$$
in \eqref{difsolu}. Consequently,
$$\int_{\zeta_0}^{\zeta_1} \omega_{\lambda,f}=0$$
for any $f$ in the range of $\lambda-A$. The same argument as in Case \uppercase\expandafter{\romannumeral2} then shows that $\lambda-A$ is not surjective, hence $\lambda\in\sigma(A)$. Therefore, we have 
$$\sigma(A)=\lfloor-\infty, \gamma_2\rceil\cup\lfloor\gamma_0, \gamma_1\rceil$$
in this case.

 


\noindent\emph{ Case \uppercase\expandafter{\romannumeral4}: $\gamma_0\leq\gamma_2$.} As shown in Case \uppercase\expandafter{\romannumeral2}  and Case \uppercase\expandafter{\romannumeral3} , both $\lfloor-\infty,\gamma_2\rceil$ and $\lfloor\gamma_0,\gamma_1\rceil$ are contained in $\sigma(A)$, hence by Corollary \ref{spectralbound}, 
$$\sigma(A)=\lfloor-\infty,\gamma_1\rceil.$$



\noindent\emph{ Case \uppercase\expandafter{\romannumeral5}.} The only remaining situation is when $\gamma_0=\gamma_1>\gamma_2$. As shown in Case \uppercase\expandafter{\romannumeral2} we have $\lfloor-\infty,\gamma_2\rceil\subset\sigma(A)$. Hence it suffices to prove that the line $\{\lambda\in\mathbb{C} : \operatorname{Re}\lambda=\gamma_0 \}$ is contained in $\sigma(A)$. 

Assume, for contradiction, that there exists some $\lambda$ on this line with $\lambda\notin\sigma(A)$. Then $\lambda-A$ is injective, which means that $\frac{\mathrm{e}^{\lambda h(z)}}{v(z)}$ does not belong to $A^p$. However, the proof of Theorem \ref{eigenvalue} shows that $\frac{\mathrm{e}^{\lambda h(z)}}{v(z)}\in A^p_\zeta$ for every $\zeta\in \partial \mathbb{D}\backslash\{\zeta_0, \zeta_1\}$. So either  $\frac{\mathrm{e}^{\lambda h(z)}}{v(z)}\notin A^p_{\zeta_0}$ or $\frac{\mathrm{e}^{\lambda h(z)}}{v(z)}\notin A^p_{\zeta_1}$. We will show that both possibilities lead to a contradiction.

First assume that  $\frac{\mathrm{e}^{\lambda h(z)}}{v(z)}\notin A^p_{\zeta_0}$. Set $\tilde{v}(z)=\frac{v(z)}{(z-\zeta_1)^c}$ and $\tilde{u}_t=\frac{\tilde{v}\circ\varphi_t}{\tilde{v}}$. Choose the constant $c$ so that $\gamma_2<\tilde{\gamma_1}<\gamma_0$. Then by Case \uppercase\expandafter{\romannumeral2} we have $\lambda\in\sigma(\tilde{A})$, where $\tilde{A}$ denotes the generator of $(\tilde{u}_tC_{\varphi_t})$. Since $v$ and $\tilde{v}$ are equivalent near $\zeta_0$, we have $\frac{\mathrm{e}^{\lambda h(z)}}{v(z)}\notin A^p_{\zeta_0}$, thus $\lambda-\tilde{A}$ is injective and consequently $\lambda-\tilde{A}$ is not surjective. On the other hand, since $\lambda-A$ is surjective, for every $f\in A^p$ there exists a constant $K_0\in\mathbb{C}$ such that
$$\frac{ \mathrm{e}^{\lambda h(z)}}{v(z)}\left(K_0-\int_0^z \omega_{\lambda,f}\right)$$
belongs to $A^p$, and in particular it belongs to $A^p_{\zeta_0}$.  The equivalence of $v$ and $\tilde{v}$ near $\zeta_0$ then implies
$$\frac{ \mathrm{e}^{\lambda h(z)}}{\tilde{v}(z)}\left(K-\int_0^z \tilde{\omega}_{\lambda,f}\right)\in A_{\zeta_0}^p$$
for some $K\in\mathbb{C}$, where 
$$\tilde{\omega}_{\lambda,f}=\mathrm{e}^{-\lambda h(\xi)}h'(\xi)\tilde{v}(\xi)f(\xi) \mathrm{d}\xi.$$
Applying Claims \ref{others} and \ref{>j} we deduce that
$$\frac{ \mathrm{e}^{\lambda h(z)}}{\tilde{v}(z)}\left(K-\int_0^z \tilde{\omega}_{\lambda,f}\right)\in A^p.$$
But this means that $\lambda-\tilde{A}$ is surjective, which is a contradiction.

If instead we have $\frac{\mathrm{e}^{\lambda h(z)}}{v(z)}\notin A^p_{\zeta_1}$, the argument is essentially the same: simply define $\tilde{v}(z)=\frac{v(z)}{(z-\zeta_0)^c}$ with a suitable $c$ so that $\gamma_2<\gamma_1<\tilde{\gamma}_0$ and then repeat the reasoning.

Combining all five cases we complete the proof of Theorem \ref{main}.
\end{proof}

\end{document}
